% titel: Blatt 0, A8: Lineares Gleichungssystem

% Teil I: Werkzeuge, kein Lösungsweg.

% Grundidee: Zwei Geraden, ein Schnittpunkt. Das Verfahren eliminiert eine
% Variable, sodass man eine einfache Gleichung erhält, dann die andere.

Zwei lineare Gleichungen in $x$ und $y$ sind geometrisch zwei Geraden; ihre
Lösung ist der Schnittpunkt. Das \emph{Additionsverfahren} und das
\emph{Einsetzungsverfahren} sind zwei Wege, diesen Schnittpunkt zu
berechnen, indem man die Unbekannte eliminiert.

\begin{werkzeugkasten}[Additionsverfahren: die Idee]
Die Grundidee ist, \emph{eine Variable wegfallen zu lassen}. Man
multipliziert eine Gleichung so mit einem Faktor, dass bei der Addition
der beiden Gleichungen die eine Variable verschwindet:
\begin{equation}
  -3x+4y=-1,\qquad 2x-5y=-4
\end{equation}
Multipliziert man die erste mit $2$ und die zweite mit $3$, steht in beiden
bzw. zusammen ein ganzer Vielfachen-Ausgleich, sodass sich $x$ beim Addieren
aufhebt ($-6x+6x=0$). Übrig bleibt eine lineare Gleichung in $y$.

\textbf{Warum funktioniert das ``Faktor-Multiplizieren``?} Eine
Gleichung ändert sich nicht, wenn man beide Seiten mit derselben Zahl
multipliziert --- die Lösungsmenge bleibt identisch. Das ist die
Harmlosigkeit der Operation.
\end{werkzeugkasten}

\begin{werkzeugkasten}[Einsetzungsverfahren: die Idee]
Eine Gleichung nach einer Variablen auflösen (z.\,B. nach $y$), diesen
Ausdruck in die \emph{andere} Gleichung einsetzen --- damit ist die
eine Variable eliminiert, man erhält eine Gleichung in nur noch einer
Unbekannten; lösen und zurücksetzen.
\end{werkzeugkasten}

\begin{achtung}[Typische Fehler]
\begin{itemize}
  \item \textbf{Koeffizienten und Vorzeichen} korrekt übernehmen
        ($-3x$, $4y$, $2x$, $-5y$) --- ein falsches Minus beim
        ``Faktor-Multiplizieren`` verschiebt den Schnittpunkt.
  \item \textbf{Beide Seiten} der Gleichung mit dem Faktor multiplizieren,
        nicht nur die linke (sonst ist die Gleichung nicht mehr
        äquivalent).
  \item \textbf{Rückprobe in beiden} Gleichungen: $(x,y)$ muss \emph{beide}
        erfüllen, nicht nur eine.
  \item Beim Einsetzungsverfahren den Ausdruck sauber klammern, damit beim
        Einsetzen keine Vorzeichen- und Rechenfehler entstehen.
\end{itemize}
\end{achtung}
